\section{Experimental design}
\label{sec:design}

\subsection{Study A: broad Transformer tensor}

The exploratory tensor contains \PhaseTensorRuns{} trained small decoder models
across \PhaseTensorConditions{} conditions and five full-pipeline seeds. It
varies feed-forward type, residual and normalization choices, objective
interpolation, optimizer geometry, width/depth/context/data scaling paths, and
synthetic-to-natural data carriers. The screen records train, IID, and OOD risk;
semantic task order; attention and MLP interventions; gradient and update
geometry; activation occupancy and spectral diagnostics; trajectory statistics;
and resource accounting. Its purpose is breadth and failure discovery.

\subsection{Study B: frozen Transformer confirmations}

The confirmation suite contains \PhaseConfirmRuns{} models in
\PhaseConfirmConditions{} conditions with \PhaseConfirmSeeds{} fresh seeds per
condition. Four feed-forward variants are compared with matched parameter
budgets where applicable. Optimizer learning rates are selected on calibration
slices and frozen for evaluation. Data-scaling paths remove the cap artifact
found in the screen. Natural-corpus byte ranges are non-overlapping and report
bits per byte together with train--test and cross-corpus diagnostics.

\subsection{Study C: portable GPU reference grid}

The public reference study is intentionally small enough to audit at raw-seed
level: widths $d\in\{16,32,64\}$, requested sample coefficients
$\alpha\in\{0.5,2,8\}$, and \ReferenceSeeds{} independent full-pipeline seeds,
for \ReferenceRuns{} trained models. Each one-layer pre-normalized Transformer
uses four heads, a GELU feed-forward block with ratio 2, AdamW, 24 update steps,
and a controlled retrieval task. It records 81 aggregate fields, including
separate risks, causal ablations, gradient geometry, attention entropy,
susceptibility, Binder cumulant, and raw outer-seed samples.

An audit found that the input pipeline enforces at least 32 training examples.
Consequently, $(d,\alpha)=(16,0.5)$ and $(16,2)$ both use 32 examples and are
the same effective control point. They remain visible as a duplicated point and
are not counted as two independent locations on a response curve. This finding
motivates the dense confirmatory design in Appendix~\ref{app:next-experiment}.

\subsection{Study D: cross-domain and continuation stress tests}

Controlled diffusion, reinforcement-learning, and multi-agent simulators vary
reverse time/locality, KL pressure/verifier noise, and field/coupling,
respectively. Their role is to test whether the PhaseCard can express
trajectory, feedback, and population observables without reusing Transformer
semantics. The predictive-continuation benchmark contains many simulator tasks
and nested replicas, but because the response model is generated by the same
controlled family, it validates holdout and censoring mechanics only.

\subsection{Compute and portability}

All trainable runs are specified independently of site paths. Repository, data,
manifest, interpreter, and output locations are runtime inputs. GPU use is an
execution coordinate, not part of the scientific identity. Public artifacts do
not contain hostnames, scheduler partitions, user directories, or private data
locations.

\FloatBarrier
